\usepackage[utf8]{inputenc}
\usepackage[english]{babel}
\usepackage[T1]{fontenc}
\usepackage{amsmath}
\usepackage{amsfonts}
\usepackage{amssymb}
\usepackage{amsthm}
\usepackage{tikz}
\usepackage{enumitem}
\usetikzlibrary{calc, intersections, fpu}
\usetikzlibrary{arrows.meta} %cool arrows
\usetikzlibrary{decorations.text}
\usepackage{tikz-cd}
\usepackage[draft]{hyperref}
\usepackage{comment}
\usepackage{geometry}
\usepackage{dsfont}
\usepackage{cleveref}


%%%%%%%%%%%%%%%%%%%%%%%%%%%%%%%%%%%%%%%%%%%%%%%%%%%%%%%%%%%%%%%%%%%%%%%%%%%%%%%%%%%
%%%%%%%%%%% Maccros/Commandes
\newcommand{\R}{\mathbb{R}}
\newcommand{\To}{\mathbb{T}}
\newcommand{\Sp}{\mathbb{S}}
\newcommand{\T}{\mathcal{T}}

%tikz command, simplifie l'écriture des neouds
\newcommand{\control}[2]{ ($#2+#1$) .. #1}
\newcommand\tikzmark[5][0]{\tikz[overlay,remember picture,baseline] \node [rotate=#1,anchor=base,xshift=#4,yshift=#5] (#2) {$\scriptstyle#3$};}
\newcommand\modulo[3]{\pgfmathparse{mod (#1,#2)}\pgfmathtruncatemacro{#3}{\pgfmathresult}}
\newcommand\euclidiv[3]{\pgfmathparse{int(#1/#2)}\pgfmathtruncatemacro{#3}{\pgfmathresult}}

%%%%%%Specific to this one
\newcommand{\cross}{\operatorname{cr}}
\newcommand{\addcross}{\operatorname{Top}}
\newcommand{\add}{\operatorname{Add}}
\newcommand{\unlink}{\operatorname{CC}}

%%%%%Environments
\newtheorem{theorem}{Theorem}[section]
\newtheorem{definition}[theorem]{Definition}
\newtheorem{proposition}[theorem]{Proposition}
\newtheorem{corollary}[theorem]{Corollary}
\newtheorem{remark}[theorem]{Remark}
\newtheorem{lemma}[theorem]{Lemma}
\newtheorem{conjecture}[theorem]{Conjecture}

%%%%%%%%%%%%%%%%%%%%%%%%%%%%%%%%%%%%%%%%%%%%%%%%%%%%%%%%%%%%%%%%%%%%%%%%%%%
